\documentclass[11pt]{article} % Times New Roman, 12pt

% SINGLE OR DOUBLE SPACE
%\usepackage{gscale_thesis_singlespace} % Single spaced thesis
\usepackage{../settings/template/gscale_thesis_doublespace} % Double spaced thesis

\usepackage{../settings/template/fancyheadings} % Header and footer styling

\usepackage[]{natbib} % Bibliography

\usepackage{setspace} % Allows double spacing but skips headers/footers
\setcounter{tocdepth}{1} % Limits the TOC to chapter and section names

% Additional packages
\usepackage{graphicx} % Allows the inclusion of figures
\usepackage{subcaption} % Allows captions to be added to subfigures
\usepackage[justification=centering]{caption} % Centres caption text
\usepackage{array} % Used for table formatting
\newcolumntype{P}[1]{>{\raggedright\let\newline\\\arraybackslash\hspace{0pt}}m{#1}}
\usepackage{booktabs} % Fancy-style tables
\usepackage{longtable} % Allows for tables that are more than one page long
\usepackage{float} % Better figure placement control
\usepackage{rotating}
\usepackage{enumerate} % Numbered lists
\usepackage[shortlabels]{enumitem} % For controlling enumerate labels
\usepackage[shortcuts]{extdash} % Allows manual hyphenation of hypenated words
\usepackage{amsmath} % Non-standard math symbols
\usepackage{amsfonts} % Extended fonts for mathematics
\usepackage{multirow}

\usepackage[hidelinks]{hyperref} % Linking to LaTeX labels and external URLs

\numberwithin{equation}{section} % Numbers equations based on their section

\usepackage{tabularx}
\usepackage{pdflscape}
\usepackage{booktabs} % for better horizontal lines
\usepackage{makecell} % for line breaks inside cells

% CUSTOM SETTINGS ********************************
\input{../settings/custom/packages}
\input{../settings/custom/settings}
\newcommand{\indep}{\rotatebox[origin=c]{90}{ $\models$ }}

\title{Characterizing the Effects of Social Disparities Using Mediation Analysis}
\author{Johann Hatzius}
\date{\today}

\begin{document}

\maketitle

\textit{This is a draft of the introduction section for the dissertation I am working on as part of my coursework for the MSc in Statistics at the London School of Economics. In this dissertation, I extend the work done on a particular class of estimators developed by statisticians working within the field causal mediation analysis. These estimators are particularly well-suited for analyzing the causal effect of health disparities between social groups, but they also have useful applications beyond this context. The research and writing that has gone into this draft is entirely my own.}

\section{Introduction}

Causal inference methods in the social and biomedical sciences are often motivated by the promise that these methods can identify and isolate the "effects of causes"~\citep{holland1986}. However, we are often interested not in the total effect of an intervention but instead in the mechanisms through which an intervention induces an effect. This is because there are many situations in which an intervention --- such as a social policy or a drug --- affects an outcome in multiple, sometimes conflicting ways. A canonical example comes from~\cite{hesslow1976}, which discusses a randomized experiment that tested whether a new birth control pill increased the risk of blood clots in women. As discussed in~\cite{pearl2001}, this drug could have a direct effect on the rate of blood clots in women by causing clots to form in some women who would otherwise not have been at risk. However, the drug could also operate through a second pathway whose sign opposes the direct effect: the birth control pill could indirectly reduce the risk of developing blood clots by reducing pregnancies, which are known to lead to higher rates of blood clots in women. A simple causal inference framing asks only: does taking the drug increase the risk of blood clots for women? However, this question does not address our concerns with the drug's safety: it is entirely possible that this drug causes blood clots in some women, but because the women who take the drug get pregnant at a lower rate, there is no observable difference in blood clot rates for women who did and did not take the drug. The policy-relevant question is what the direct effect of the drug on blood clot risk is, after accounting for the indirect effect that operates through the lowered risk of pregnancy. 

Mediation analysis attempts to elucidate the mechanisms through which an intervention affects an outcome of interest. It does so by decomposing the total causal effect of an intervention on an outcome into the causal effects of distinct pathways through which the intervention affects the outcome. Informally, the total observable causal effect of an intervention on an outcome is typically decomposed into a \textbf{direct effect} that represents the unmediated effect that an intervention has on the outcome and an \textbf{indirect effect} that represents the effect that an intervention has on an outcome only through the way that it affects mediator variables that themselves affect the outcome. 

Since statisticians first became interested in formalizing mediation analysis methods in the 1990s and early 2000s, several different versions of mediational effects have been proposed in the literature, each with their attendant interpretations, required assumptions, and theoretical issues. A central issue that statisticians have contended with is the difficulty of justifying the identification assumptions that are required to identify direct and indirect effects~\citep{vanderweele2015}. Identification in classical causal inference without mediators is already challenging, and applied work relies on randomized experiment or observational designs that approximate them~\citep{imbens2015}. Popular versions of mediational effects either require assumptions that cannot be empirically tested even under experimental conditions~\citep{imaietal2010} or carefully designed multi-armed experimental designs to admit valid causal interpretations~\citep{morenobetancur2018}. 

Natural effects, which capture the effects of an intervention through a mediator compared to the effect of the intervention had the mediator been set at the level it would have "naturally" attained under no intervention, were the first class of effects proposed in the literature ~\citep{robinsgreenland1992, pearl2001, avinetal2005}. While theoretically appealing, these effects have been controversial because they require an empirically untestable "cross-world" assumption for identification, and the practical usefulness of their implications has also been called into question~\citep{vanderweelevansteelandt2009, imaietal2010}. Interventional effects, which require significantly weaker assumptions but do not capture the same notion of mediation as natural effects~\citep{miles2023}, have recently become the dominant target estimand in applied work~\citep{vanderweeleetal2014, vanderweele2015, vansteelandtdaniel2017}. 

Applied researchers have recently become interested in the applications of mediation analysis to problems in the health disparities literature. This has been driven by~\cite{vanderweelerobinson2014}, who propose a version of interventional effects that can be used as a tool to help characterize the effects of race and sex based disparities. Here, a disparity in some intermediate variable - for example, educational outcomes - is observed between different social groups and the causal question of interest is how equalizing the distribution of this outcome between these groups would affect an outcome of interest. For example: if a disparity in educational attainment is observed between black and white Americans, what would be the effect of performing some intervention on black American educational attainment to equalize it to the distribution of white American educational attainment on black American lifetime earnings? Although this version of interventional effect is unnamed by~\cite{vanderweelerobinson2014}, I will refer to this type of mediational effect as a \textbf{VR-interventional effect}. 

VR-interventional effects occupy an appealing conceptual middle ground compared to other popular estimators that have been used to try to decompose the effects of discrimination~\citep{jacksonvanderweele2018, jackson2021}. The Oaxaca-Blinder decomposition cannot be used to make claims about the causal effect of discriminatory policies or disparities and is susceptible to bias through selection and confounding. Other causal inference estimators (e.g. natural effects estimators), meanwhile, both require more stringent assumptions and problematically conceptualize immutable characteristics such as race and sex as manipulable causal characteristics. In the causal inference literature, many researchers find the idea that unchangeable characteristics like race, sex, and (in sociology) social class can have a causal interpretation conceptually dubious~\citep{imbens2015}. This norm is exemplified by a well-known slogan coined by Holland: "no causation without manipulation!"~\citep{holland1986}. A focus on VR-interventional effects avoids postulating a "causal effect of race" while focusing on the causal effect that disparities in intermediate outcomes that might be a result of existing social inequalities have on ultimate outcomes of interest. 

The main contribution of this paper is to formalize the VR-interventional effects discussed in~\cite{vanderweelerobinson2014} and ~\cite{jacksonvanderweele2018} by connecting these ideas to a separate strand of research on estimating the casual effects of population level interventions. This research originates from work done in~\cite{hubbardvanderlaan2008} and~\cite{diazvanderlaan2012}, where a framework is developed for (potentially stochastic) population level interventions whose effects depend on a known existing pre-treatment distribution of a treatment in the population. This theoretical setup has many similarities to the one developed in this paper; the key difference is that in the~\cite{hubbardvanderlaan2008} setup, the untreated and treated distributions that are being assigned are known, while in the setup developed for VR-interventional effects in this paper, these distributions are unknown and must be estimated. This extension introduces a new source of estimation uncertainty that must be accounted for, meaning that the estimators developed in the population intervention literature cannot be entirely ported over for VR-interventional effects, and requires a different set of identification assumptions. The connection to the population intervention literature also shows that VR-interventional effects have applications beyond disparities research. 

After the connection between these literatures is developed, I then extend the population intervention framework for VR-interventional effects to cases with multiple mediators. Applied researchers frequently encounter problems where there are multiple mediators of a causal effect which are also themselves causally ordered. For example,~\cite{morenobetancur2018} discuss a causal structure where headache and insomnia mediate the effect of an anti-depression treatment on depression, and insomnia itself also mediates the effect of the treatment on headache. The estimation of VR-interventional effects with multiple mediators has been briefly discussed in~\cite{jacksonvanderweele2018} and~\cite{jackson2021} and was developed for normal interventional effects by~\cite{vansteelandtdaniel2017}, but analogues have not been developed in the population intervention literature. While the extension to cases with multiple mediators is relatively straightforward in terms of identification assumptions, estimation becomes much more challenging. 

Having developed this theoretical machinery, I demonstrate an application of VR-interventional effects in a multiple mediator setting.~\cite{kuhaetal2021} looked at how ideas from mediation analysis could be used to characterize the mediating role of education in post-war British social class mobility. The definition of social class used by~\cite{kuhaetal2021} is one that, like race or sex, is seen as an immutable characteristic, which motivates the focus on the estimation of associations instead of causal effects in the original paper. However, VR-interventional effects can also be applied to study the causal effect of disparities in educational outcomes between social classes. In my analysis, the original study is extended to look at the role of measured intelligence as an additional mediator of social class mobility. Finally, as suggested by~\cite{jackson2021}, I propose a sensitivity analysis that can be used for VR-interventional effects estimation and demonstrate how to conduct this analysis.  

Section 2 of this paper outlines an in-depth historical review of the mediation analysis literature. Section 3 develops the formal framework and notation for mediation analysis and details the identification assumptions, conceptual interpretations, and estimation strategies for natural, controlled, and interventional effects. Section 4 develops the formal setting for VR-interventional effects with one mediator. Section 5 extends the framework developed in Section 4 to cases with multiple mediators. Section 6 walks through the application of VR-interventional effects in a multiple mediator setting. Section 7 concludes. 


\bibliographystyle{IEEEtranN}
\bibliography{../bib/references,../bib/secondary_references}

\end{document}

